\documentclass{beamer}

\usepackage{fontspec}
\setmainfont{CMU Serif}
\setsansfont{CMU Sans Serif}
\setmonofont{CMU Typewriter Text}

\usepackage{graphicx}

\usepackage{polyglossia}
\setmainlanguage[babelshorthands=true]{russian}

\usepackage{minted}
\usepackage[svgnames]{xcolor}
\setminted{autogobble=true, bgcolor=Beige, bgcolorpadding=0.2em}

\usepackage{subcaption}

\newcommand\comment[1]{\textcolor{magenta}{\textit{{#1}}}}

\usepackage{xspace}
\newcommand\cpp{C\texttt{++}\xspace}

\usepackage{minted}

\usepackage{hyperref}
\hypersetup{
    colorlinks=true,
    linktoc=all,
    linkcolor=black,
}
\urlstyle{same}

\usetheme{Madrid}
\usecolortheme{beaver}

\title{Линковка в C\texttt{++}}
\author{Сениченков Пётр}
\date{2026}

\begin{document}
\frame{\titlepage}

\section{Что такое линковка}

\begin{frame}{Линковка}
    \emph{Линковка} "--- это этап сборки, на котором отдельные единицы трансляции
    (\emph{translation unit}), единицы подстановки шаблонов
    (\emph{instantiation unit}) и внешние библиотеки собираются в исполняемый файл или
    библиотеку (\emph{program image}).

    \footnotetext{\url{https://github.com/p-senichenkov/cpp-linking-talk}}

    \pause
    \begin{quote}
        All external entity references are resolved.
        Library components are linked to satisfy external references to entities not defined in the
        current translation.
        All such translator output is collected into a program image which contains information
        needed for execution in its execution environment.
    \end{quote}
    \begin{quote}
        \raggedleft
        ISO/IEC 14882:2020, 5.2, translation phase 9
    \end{quote}
\end{frame}

\begin{frame}{Немного терминологии}
    Чтобы не запутаться, будем использовать не совсем стандартные термины (если не очевидно, о чём
    идёт речь):
    \begin{itemize}
        \item Compile-time
            \begin{itemize}
                \item \textbf{Компоновка}
                \item 9 фаза трансляции
                \item Link-editor
            \end{itemize}
        \item Runtime
            \begin{itemize}
                \item \textbf{Системный интерпретатор}
                \item \texttt{ld.so}
                \item Runtime linker
                \item Dynamic linker/loader
            \end{itemize}
    \end{itemize}
\end{frame}

\begin{frame}{Что делает линкер}
    \begin{itemize}
        \item Разрешение символов
        \item Объединение секций
        \item Address relocation
        \item Подключение динамических библиотек
            \pause
            \begin{itemize}
                \item Position independent code (PIC)
                \item Compile-time
                    \begin{itemize}
                        \item Создаются Global offset table (GOT) и Procedure linkage table (PLT)
                        \item Заполняются динамические секции ELF
                    \end{itemize}
                \item Runtime
                    \begin{itemize}
                        \item Запускается \texttt{ld.so <executable>}
                        \item Находит библиотеки из \texttt{DT\textunderscore NEEDED}, обходит их
                            рекурсивно и загружает
                        \item Заполняет адреса в GOT
                        \item При первом вызове заполняется PLT
                    \end{itemize}
            \end{itemize}
    \end{itemize}
\end{frame}

\section{Особенности линковки в \cpp}

\begin{frame}{Storage duration, linkage, visibility}
    \begin{itemize}
        \item Storage duration
            \begin{itemize}
                \item Static
                \item Thread
                \item Automatic
                \item Dynamic
            \end{itemize}
        \item Linkage
            \begin{itemize}
                \item No linkage
                \item Internal
                \item External
                    \begin{itemize}
                        \item Language
                    \end{itemize}
                \item Module
            \end{itemize}
        \item Visibility
    \end{itemize}
\end{frame}

\begin{frame}{Name mangling, language linkage}
    \begin{itemize}
        \item Name mangling
        \item Language linkage
            \begin{itemize}
                \item \mintinline[bgcolor=white]{C++}{extern "C++"}
                \item \mintinline[bgcolor=white]{C++}{extern "C"}
                \begin{itemize}
                    \item Отключает name mangling
                \end{itemize}
            \item Реализации могут поддерживать другие языки
                \begin{itemize}
                    \item На практике это не используется
                \end{itemize}
            \end{itemize}
    \end{itemize}
\end{frame}

\begin{frame}[fragile]{Name mangling, language linkage}
    \begin{minted}[mathescape]{C++}
        int foo(double, long);            // $ \to $ _Z3foodl
        extern "C" int bar(double, long); // $ \to $ bar

        extern "C" typedef void FUNC();
        FUNC* baz;                        // $ \to $ baz
        extern "C" FUNC* qux;             // $ \to $ qux

        extern "C" {
        struct C {
            void mem(void (*)());         // $ \to $ _ZN1C3memEPFvvE
        };
        }

        struct D {
            void mem(void (*)());         // $ \to $ _ZN1D3memEPFvvE
        };
    \end{minted}
\end{frame}

\begin{frame}{ODR}
    \begin{itemize}
        \item \textbf{Не больше одного} определения в одной translation unit
        \item \textbf{Ровно одно} определение для каждой не-\texttt{inline} функции или
            \texttt{inline}-переменной, если она используется
        \item \textbf{Полное} определение для каждого класса, если он используется
        \item Для класса, \texttt{inline}-функции, \texttt{inline}-переменной, перечисления,
            шаблона (но не полной специализации) разрешено несколько определений, если они
            \textit{полностью совпадают} и определены в разных translation unit (модулях)
    \end{itemize}
\end{frame}

\begin{frame}{Глобальные переменные и copy relocation}
    \begin{itemize}
        \item Copy relocation
            \begin{itemize}
                \item Сам исполняемый файл обычно создаётся из position-dependent кода.
                    Нельзя использовать GOT и PLT для разрешения адресов
                \item Когда компоновщик видит ссылку на данные, определённые в разделяемом объекте,
                    он создаёт пустую запись в \texttt{.bss}, того же размера, что и эти данные
                \item Интерпретатор копирует данные в эти записи
            \end{itemize}
    \end{itemize}
\end{frame}

\begin{frame}{Шаблоны и линковка}
    \begin{itemize}
        \item Borland model
            \begin{itemize}
                \item Компилятор подставляет шаблоны там, где они используются
                \item Линкер ``склеивает'' одинаковые подстановки
                \item Нужно подменять линкер
                \item Дольше компиляция
            \end{itemize}
            \pause
        \item Cfront model
            \begin{itemize}
                \item Глобальный реестр подстановок
                \item Можно использовать системный линкер
                \item Гораздо сложнее в реализации
            \end{itemize}
    \end{itemize}
\end{frame}

\begin{frame}{Виртуальные функции и key function}
    \begin{itemize}
        \item vtable размещается в объектном файле, который содержит \emph{key function}
            \begin{quote}
                Key function is defined as the first non-inline, non-pure, virtual function declared
                in the class. If you do not provide a definition for the key function (or fail to
                link to the file providing the definition) then the linker will not be able to find
                the class' vtable.
            \end{quote}
            \begin{itemize}
                \item Если \textbf{не объявлено} ни одной не-inline, не чистой, виртуальной функции,
                    то компилятор сгенерирует vtable в каждой TU, которая ссылается на
                    этот класс (в секции COMDAT)
            \end{itemize}
    \end{itemize}
\end{frame}

\begin{frame}{LTO}
    \begin{itemize}
        \item Компилятор записывает промежуточное представление (GIMPLE, LLVM bitcode) в объектные
            файлы
        \item Link-time optimizer читает промежуточное представление и работает вместе с линкером
            \begin{itemize}
                \item Линкер подключает \texttt{liblto\textunderscore plugin.so} или
                    \texttt{libLTO.so}
            \end{itemize}
        \pause
        \item В основном, те же оптимизации, что и при компиляции, но оптимизатор видит несколько
            TU:
            \begin{itemize}
                \item Cross-module inlining
                \item Dead code elimination
                \item Devirtualization
                \item Identical code folding
                \item Code layout optimization
                \item \texttt{COMDAT} folding
            \end{itemize}
    \end{itemize}
\end{frame}

\section{Линковка в Rust}

\begin{frame}{Линковка в Rust}
    \begin{itemize}
        \item Единица компиляции "--- crate
        \item Экспорт символов через модули, а не текстовые подстановки
        \item Name mangling
            \begin{itemize}
                \item Содержит имя и хэш crate
            \end{itemize}
        \item Нет стабильного ABI
        \item \texttt{rlib} содержит дополнительные метаданные
            \begin{itemize}
                \item Версия компилятора
                \item Crate id
                \item Информация об исходных файлах
                \item Информация об экспортированных символах
                \item Mid-level IR
                \item LLVM bitcode
            \end{itemize}
        \item Generic monomorphization
    \end{itemize}
\end{frame}

\begin{frame}[fragile]{Rust foreign function interface}
    \begin{itemize}
        \item C ABI
        \item unsafe
            \begin{itemize}
                \item Линковка в runtime
                \item Коллизии имён без mangling
            \end{itemize}
    \end{itemize}

    \begin{minted}{rust}
        #[link(name = "my_c_library")]
        unsafe extern "C" {
            fn my_c_function(x: i32) -> bool;
        }
    \end{minted}

    \begin{minted}{rust}
        #[unsafe(no_mangle)]
        pub extern "C" fn callable_from_c(x: i32) -> bool {
            x % 3 == 0
        }
    \end{minted}
\end{frame}

\section{Линковка в go}

\begin{frame}{Линковка в Go}
    \begin{itemize}
        \item Единица сборки "--- package
            \begin{itemize}
                \item Export data и метаданные
            \end{itemize}
        \item В основном используется внутренний линкер
        \item Self-contained binary
            \begin{itemize}
                \item Динамические библиотеки практически не используются
            \end{itemize}
        \item Относительно стабильный ABI
    \end{itemize}
\end{frame}

\begin{frame}[fragile]{Inter-language linking in Go}
    \begin{minted}{go}
        package pkg

        // #cgo LDFLAGS: -lpthread
        // #include <stdio.h>
        // void myprintf(char* s) { printf("%s\n", s); }
        import "C"
        import "unsafe"

        //export MyFunction
        func MyFunction(arg1, arg2 int) int64 {...}

        func main() {
            cs := C.CString("Hellow from stdio")
            C.myprint(cs)
            C.free(unsafe.Pointer(cs))
        }
    \end{minted}
\end{frame}

\begin{frame}[fragile]{Как это выглядит в коде}
    \begin{itemize}
        \item Псевдо-пакет \texttt{"C"}
        \item Комментарий перед \texttt{import "C"} содержит
            \begin{itemize}
                \item Флаги для компиляторов
                    (\mintinline[bgcolor=white]{go}{// #cgo <environmental variable>})
                \item Общий заголовок для сгенерированного кода на C
            \end{itemize}
        \item Экспорт функций с помощью комментария \mintinline[bgcolor=white]{go}{//export <name>}
        \item Обращение к переменным и функциям C через \mintinline[bgcolor=white]{go}{C.<name>}
        \item Не все возможности C поддерживаются
            \begin{itemize}
                \item Variadic функции
                \item Вызов функции по указателю
            \end{itemize}
        \item Нужно постоянно приводить типы
        \item Хрупкая работа с указателями
    \end{itemize}
\end{frame}

\begin{frame}{Как это работает}
    \begin{itemize}
        \item Вызывается не компилятор, а \texttt{go tool cgo}
        \item Собирает все исходники C, \cpp и Fortran в текущей директории
            \begin{itemize}
                \item Но не в поддиректориях
            \end{itemize}
        \item Генерирует несколько файлов на C и Go и собирает их
            \begin{itemize}
                \item Экспортируемые файлы на C для подключения из внешних программ
                \item Обёртки над использованными функциями на C
                \item Исходный файл на Go с подставленными типами и именами из псевдо-модуля
                    \texttt{"C"}
                \item Вспомогательные файлы с обёртками над типами, некоторыми статическими
                    проверками, etc.
                \item Подробнее "--- в репозитории с примерами
            \end{itemize}
    \end{itemize}
\end{frame}

\section{Системы сборки}

\begin{frame}[fragile]{Системы конфигурации и сборки для \cpp}
    \begin{itemize}
        \item Системы сборки
            \begin{itemize}
                \item Просто запускают команды
                \item topsort
            \end{itemize}
        \item Системы конфигурации
            \begin{itemize}
                \item Мало знают о структуре программы
            \end{itemize}
    \end{itemize}

    \pause
    \begin{minted}{cpp}
        // main.cpp
        #include "lib.h"
    \end{minted}
    \begin{minted}{cmake}
        # CMakeLists.txt
        add_executable(MyProgram main.cpp)
        target_link_libraries(MyProgram MyLibrary)
    \end{minted}
\end{frame}

\begin{frame}[fragile]{Системы сборки для Rust и go}
    \begin{itemize}
        \item Cargo
            \begin{itemize}
                \item Package manager
                \item Build orchestrator
                \item Структура проекта стандартизирована
                    \begin{itemize}
                        \item \texttt{Cargo.toml}
                        \item \texttt{src/main.rs} или \texttt{src/lib.rs}
                    \end{itemize}
            \end{itemize}
        \item Go toolchain
            \begin{itemize}
                \item Единый toolchain, тесно интегрированный в сам язык
                \item Build graph $ \leftrightsquigarrow $ import graph
                \item Разработчик практически не контролирует сборку
            \end{itemize}
    \end{itemize}
\end{frame}

\end{document}
